{\scriptsize
\setlength{\tabcolsep}{3pt}
\begin{longtable}{>{\raggedright\arraybackslash}p{2.19cm}>{\raggedright\arraybackslash}p{2.19cm}>{\raggedright\arraybackslash}p{2.19cm}>{\raggedright\arraybackslash}p{2.19cm}>{\raggedright\arraybackslash}p{2.19cm}>{\raggedright\arraybackslash}p{2.19cm}>{\raggedright\arraybackslash}p{2.19cm}}
\caption{Nombre de composantes actives sur 20 initialisations indépendantes}\label{tab:2_11}\\
\toprule
\textbf{initialisation} & \textbf{algorithme} & \textbf{n\_repetitions} & \textbf{3 actives} & \textbf{4 actives} & \textbf{5 actives} & \textbf{6 actives} \\
\midrule
\endfirsthead
\toprule
\textbf{initialisation} & \textbf{algorithme} & \textbf{n\_repetitions} & \textbf{3 actives} & \textbf{4 actives} & \textbf{5 actives} & \textbf{6 actives} \\
\midrule
\endhead
\bottomrule
\endfoot
init k-means & EM classique & 20 & 0 & 0 & 0 & 20 \\
init k-means & EM variationnel (priors diffus) & 20 & 0 & 12 & 8 & 0 \\
init k-means & EM variationnel (priors empiriques) & 20 & 20 & 0 & 0 & 0 \\
init aléatoire & EM classique & 20 & 0 & 0 & 0 & 20 \\
init aléatoire & EM variationnel (priors diffus) & 20 & 4 & 14 & 2 & 0 \\
init aléatoire & EM variationnel (priors empiriques) & 20 & 20 & 0 & 0 & 0 \\
\end{longtable}
}
